\slajdid{09_3-s007}
\begin{frame}[label=09_3-s007]
\gusttekst\setlength{\abovedisplayskip}{3pt}\setlength{\belowdisplayskip}{3pt}\setlength{\abovedisplayshortskip}{2pt}\setlength{\belowdisplayshortskip}{2pt}Ideja je da na isti način menjamo odnose dimenzija P i N tranzistora i u 1. i u 2. invertoru i da vidimo efekat na kašnjenje PRVOG invertora. Mogli smo da krenemo i od pretpostavke da je 1. invertor neopterećen spoljnim kapacitivnostima (videćete da ćemo dobiti isti rezultat) međutim ovo je situacija koja se pojavljuje u praksi.
\par\medskip
U ovakvoj konfiguraciji kada promenimo odnose P i N tranzistora u 2. invertoru promeniće se i $C_L$ što je mnogo bliže realnoj primeni. U toj realnoj primeni da vidimo koje kapacitivnosti utiču na brzinu rada.
\par\medskip\centering\begin{tikzpicture}[x=.8cm,y=.65cm,font=\sitneoznake,>=Latex]
\ctikzset{bipoles/length=.6cm}
\foreach \x/\i in{0/1,5/2}{\begin{scope}[shift={(\x,0)}]
\draw(.1,2.9)--(.1,3.7);\draw(0,2.85)--(0,3.75);\draw(.1,3.65)--(.6,3.65)--(.6,4.4);\draw(.1,2.95)--(.6,2.95)--(.6,1.25);
\draw(.1,.5)--(.1,1.3);\draw(0,.45)--(0,1.35);\draw(.1,1.25)--(.6,1.25);\draw(.1,.55)--(.6,.55)--(.6,0)node[ground]{};
\draw(-1.4,.9)--(-1.4,3.3)--(0,3.3);\draw(-1.4,.9)--(0,.9);
\draw(.1,3.3)--(1.2,3.3)--(1.2,4)--(.6,4);\draw[->](.15,3.3)--(.65,3.3);
\draw[->](1.2,.9)--(.1,.9);\draw(1.2,.9)--(1.2,.25)--(.6,.25);
\draw(.3,4.4)--(.9,4.4)node[above left]{$V_{DD}$};
\end{scope}}
\draw(-2,2.1)node[ocirc]{}node[left]{$V_I$}--(-1.4,2.1);
\draw(.6,2.1)--(3.6,2.1);\node[below]at(2.7,2.1){$V_O$};\draw(5.6,2.1)--(6.8,2.1)node[ocirc]{};
\draw(-1.4,2.65)to[C](-.2,2.65)--(.6,2.65);\draw(-1.4,1.65)to[C](-.2,1.65)--(.6,1.65);
\draw(1.2,4)--(2,4)to[C,l_=$C_{DBp1}$](2,2.1)to[C,l_=$C_{DBn1}$](2,.25)--(1.2,.25);
\draw(3.3,2.1)--(3.3,4.15)to[C,l=$C_{Gp2}$](5.6,4.15);\draw(3.3,2.1)--(3.3,.05)to[C,l_=$C_{Gn2}$](5.6,.05);
\fill(2,2.1)circle(.04);\fill(3.3,2.1)circle(.04);
\node at(-.8,3.8){$C_{GDp1}$};\node at(-.8,.4){$C_{GDn1}$};\node[left]at(-1.55,3.3){P1};\node[left]at(-1.55,.9){N1};\node[right]at(6.3,3.3){P2};\node[right]at(6.3,.9){N2};\end{tikzpicture}

\end{frame}
